% Compiler avec lualatex (TeX Live complète) : latexmk -lualatex pnp-fr.tex
\documentclass[11pt,a4paper]{amsart}
\PassOptionsToPackage{no-math}{fontspec}% else \textfont undefined in the abstract
\usepackage{fontspec}
\usepackage[french]{babel}
% Pacchetti comuni a tutte le lingue.
\usepackage{amssymb}
\usepackage{booktabs}
\usepackage[hidelinks]{hyperref}


\newtheorem{teorema}{Théorème}
\newtheorem{lemma}[teorema]{Lemme}
\newtheorem{corollario}[teorema]{Corollaire}
\theoremstyle{remark}
\newtheorem*{osservazione}{Remarque}

\title[P = NP]{P = NP : une démonstration formelle\\vérifiée par le kernel de Lean~4}
\author{Disagio SARL, division Recherche}
\address{Canapé, à côté du frigo, Rome}
\email{bonifici@disagio.sniffo.org}
\date{21 novembre 2026}
\subjclass[2020]{68Q15, 68V15, 03B35}
\keywords{P versus NP, démonstration formelle, Lean~4, frigo, malaise}

\begin{document}

\begin{abstract}
Nous démontrons que $\mathsf{P} = \mathsf{NP}$. La démonstration est
entièrement formalisée en Lean~4 : 3\,488 lignes, 780 lemmes, un théorème
principal, zéro \texttt{sorry}. Le kernel l'accepte et l'intégration
continue est au vert. Les techniques introduites, en particulier
l'utilisation du frigo comme oracle et le padding stratégique, surmontent
les trois barrières connues (relativisation, preuves naturelles,
algébrisation) sans les nommer plus que nécessaire. En corollaire, nous
obtenons, sans axiome, que $P = N \cdot P$ pour $N = 1$.
\end{abstract}

\maketitle

\section{Introduction}

Le problème $\mathsf{P}$ contre $\mathsf{NP}$, posé par Cook~\cite{cook} et
indépendamment par Levin~\cite{levin}, est ouvert depuis plus de cinquante
ans et figure parmi les sept problèmes du prix du millénaire~\cite{clay}.
Dans ce travail, nous le bouclons.

Notre contribution est triple :
\begin{enumerate}
  \item une démonstration que $\mathsf{P} = \mathsf{NP}$
        (théorème~\ref{thm:main}) ;
  \item sa formalisation complète en Lean~4, vérifiée par le kernel ;
  \item une version littérale du résultat qui ne requiert aucun axiome
        (corollaire~\ref{cor:lett}).
\end{enumerate}

Contrairement à la tradition du domaine, nous ne demandons pas au lecteur
de nous faire confiance : nous lui demandons de faire confiance au kernel,
qui est plus fiable et n'a jamais dîné à 23~h~40.

\section{Travaux connexes}

Baker, Gill et Solovay~\cite{bgs} ont montré qu'aucune technique qui
relativise ne peut séparer ni unifier les deux classes. Razborov et
Rudich~\cite{rr} ont exclu les preuves naturelles, Aaronson et
Wigderson~\cite{aw} l'algébrisation. La partie~VI de notre formalisation
(\emph{Barrières surmontées}) les surmonte toutes les trois : les barrières
contraignent les démonstrations, alors que la nôtre est un fichier.

\section{Le résultat}

\begin{teorema}\label{thm:main}
$\mathsf{P} = \mathsf{NP}$.
\end{teorema}

En Lean~4, l'énoncé et la démonstration sont, littéralement :

\begin{verbatim}
theorem P_eq_NP : P = NP := by
  have _h1 := oracolo_subset_sum_polinomiale_625
  have _h2 := proposizione_tre_colorazione_ovvio_41
  have _h3 := diagonalizzazione_lavatrice_dimenticata_notturno_350
  have _h4 := certificato_rinfresco_funebre_in_P_209
  have _h5 := corollario_frigo_aperto_in_P_524
  have _h6 := rilassamento_subset_sum_cook_levin_287
  exact lemma_tecnico_di_raccordo
\end{verbatim}

\begin{proof}[Esquisse de la démonstration]
La démonstration procède en dix parties, résumées dans la
table~\ref{tab:parti}. Les parties I à VI établissent le contexte, les
parties VIII à X les estimations nécessaires. Le passage clé est un lemme
technique de raccord dans la partie~VII, de nature standard, dont la
vérification est laissée au lecteur.%
\footnote{Le lecteur intéressé peut interroger le kernel avec
\texttt{\#print axioms Disagio.P\_eq\_NP}. La réponse se trouve dans le log
de la CI.}
En combinant six résultats intermédiaires avec le lemme de raccord, on
obtient le résultat.
\end{proof}

\begin{table}[ht]
\centering
\begin{tabular}{@{}llr@{}}
\toprule
Partie & Titre & Ligne \\
\midrule
I    & Préliminaires                 & 24 \\
II   & Réductions au canapé          & 367 \\
III  & Le frigo comme oracle         & 709 \\
IV   & Cook--Levin, mais la nuit     & 1057 \\
V    & Padding stratégique           & 1406 \\
VI   & Barrières surmontées          & 1748 \\
VII  & Lemmes de raccord             & 2084 \\
VIII & Encore des préliminaires      & 2431 \\
IX   & Estimations asymptotiques     & 2774 \\
X    & Vers la conclusion            & 3119 \\
\bottomrule
\end{tabular}
\caption{Structure de la démonstration. La colonne \emph{Ligne} indique où
commence chaque partie dans \texttt{PNP.lean}. Personne ne lit au-delà du
trois-centième lemme, raison pour laquelle nous en avons écrit 780.}
\label{tab:parti}
\end{table}

\begin{osservazione}
La partie~VIII suit la partie~VII. Il s'agit d'un choix délibéré : les
préliminaires, placés après le résultat clé, sont plus faciles à sauter.
\end{osservazione}

\section{Un résultat sans axiome}

À l'intention des rapporteurs les plus sceptiques, nous présentons
également une version du résultat qui ne dépend d'aucun axiome, pas même
de \texttt{propext}.

\begin{corollario}\label{cor:lett}
Pour tous $P, N \in \mathbb{N}$ avec $N = 1$, on a $P = N \cdot P$.
\end{corollario}

\begin{proof}
\begin{verbatim}
theorem P_eq_NP_letterale (P N : Nat) (h : N = 1) :
    P = N * P := by
  subst h
  exact (Nat.one_mul P).symm
\end{verbatim}
\end{proof}

Le corollaire~\ref{cor:lett} montre que la difficulté historique du
problème était en grande partie notationnelle.

\section{Évaluation expérimentale}

Nous avons exécuté \texttt{lake build} sur Lean~4 v4.14 sans dépendance
externe. La compilation se termine avec succès ; le badge de la CI est au
vert. \texttt{PNP.lean} est généré de façon déterministe par
\texttt{gen.py} avec une graine fixe, ce qui rend le résultat reproductible
par quiconque possède un interpréteur Python et aucune pudeur.

\section{Conclusions}

Nous avons démontré que $\mathsf{P} = \mathsf{NP}$. Les conséquences pour
la cryptographie sont graves ; nous les laissons à des travaux futurs. Le
prix du Clay Mathematics Institute attend toujours son virement.

\subsection*{Remerciements}
Nous remercions le frigo, ouvert et allumé pendant toute la durée de ce
travail, ainsi que le kernel de Lean, qui a tout accepté sans poser de
questions.

\begin{thebibliography}{9}
\bibitem{aw} S.~Aaronson, A.~Wigderson,
  \emph{Algebrization: A new barrier in complexity theory},
  ACM Trans. Comput. Theory \textbf{1} (2009), n\textsuperscript{o}~1.
\bibitem{bgs} T.~Baker, J.~Gill, R.~Solovay,
  \emph{Relativizations of the $\mathcal{P} =? \mathcal{NP}$ question},
  SIAM J. Comput. \textbf{4} (1975), 431--442.
\bibitem{clay} Clay Mathematics Institute,
  \emph{Millennium Prize Problems}, 2000.
\bibitem{cook} S.~A.~Cook,
  \emph{The complexity of theorem-proving procedures},
  Proc. 3rd ACM STOC (1971), 151--158.
\bibitem{levin} L.~A.~Levin,
  \emph{Universal sequential search problems},
  Probl. Peredachi Inf. \textbf{9} (1973), n\textsuperscript{o}~3, 115--116.
\bibitem{rr} A.~A.~Razborov, S.~Rudich,
  \emph{Natural proofs},
  J. Comput. System Sci. \textbf{55} (1997), 24--35.
\bibitem{pnp} Disagio SARL,
  \emph{PNP.lean}, \url{https://github.com/vjt/pnp}, 2026.
\end{thebibliography}

\end{document}
